\chapter*{ВВЕДЕНИЕ}
\addcontentsline{toc}{chapter}{ВВЕДЕНИЕ}

Двигатель является ядром любого транспортного средства. Как и любой другой узел автомобиля, он должен оставаться компактным, лёгким и мощным. По этой причине двигатели проектируются высокооборотистыми, с относительно небольшим крутящим моментом. Для передачи мощности от двигателя к колёсам, преобразования момента и оборотов, а также для поддержания работы двигателя в оптимальном диапазоне во всех режимах движения применяются трансмиссии.

Трансмиссии бывают двух типов: ступенчатые и бесступенчатые. Ступенчатые трансмиссии - механические, автоматические и роботизированные - используют зубчатые передачи, где каждая пара шестерён отвечает за определённую передачу. Это ограничивает общее количество передаточных чисел, фиксирует оптимальные рабочие точки для каждой ступени и требует кратковременного разрыва потока мощности при переключении.

В свою очередь, бесступенчатые трансмиссии - вариаторы - обеспечивают плавное и непрерывное изменение передаточного отношения. Благодаря этому двигатель способен постоянно работать в наиболее эффективном диапазоне оборотов, без потерь мощности и рывков, характерных для переключений передач. Такая особенность делает вариатор особенно привлекательным с точки зрения повышения комфорта и топливной экономичности.

Целью данной курсовой работы является разработка бесступенчатой коробки передач, изучение её конструктивных особенностей, а также анализ преимуществ и недостатков по сравнению с традиционными ступенчатыми трансмиссиями.